% ===================== RESULTADOS =====================
\section{Resultados}
\label{sec:resultados}

[Presentar los resultados del estudio de forma organizada y cuantitativa.
Los resultados deben ser presentados sin interpretacion; el analisis se reserva
para la seccion posterior. Incluir tablas consolidadas, perfiles numericos
y figuras representativas.]

\begin{longtable}{@{}clcccccc@{}}
\caption{[Leyenda descriptiva de la tabla de resultados consolidados.]}
\label{tab:resumen_resultados} \\
\toprule
\textbf{Esc.} & \textbf{Descripcion} & \textbf{[Var 1]} & \textbf{[Var 2]} & \textbf{[Var 3]} & \textbf{[Var 4]} & \textbf{[Var 5]} & \textbf{Estado} \\
\midrule
\endfirsthead
\multicolumn{8}{c}{{\bfseries Tabla \thetable\ -- continuacion}} \\
\toprule
\textbf{Esc.} & \textbf{Descripcion} & \textbf{[Var 1]} & \textbf{[Var 2]} & \textbf{[Var 3]} & \textbf{[Var 4]} & \textbf{[Var 5]} & \textbf{Estado} \\
\midrule
\endhead
\bottomrule
\endfoot
1 & [Descripcion] & [Valor] & [Valor] & [Valor] & [Valor] & [Valor] & [Estado] \\
2 & [Descripcion] & [Valor] & [Valor] & [Valor] & [Valor] & [Valor] & [Estado] \\
\end{longtable}

\vspace{4pt}
\noindent\small\textit{Nota:} [Incluir notas al pie explicando simbolos, criterios de evaluacion,
referencias a metodologia, y cualquier aclaracion necesaria para interpretar la tabla.]

% ─────────────────────────────────────────────────────────────────────────────
\subsection{[Subseccion de Resultados Adicionales]}

[Presentar resultados complementarios: analisis termico, balances de masa/energia,
verificaciones independientes, etc. Cada subseccion debe tener su propia tabla
o figura con leyenda descriptiva.]

\begin{table}[htbp]
    \centering
    \caption{[Leyenda descriptiva de resultados complementarios.]}
    \label{tab:resultados_add}
    \begin{tabular}{@{}lcccc@{}}
        \toprule
        \textbf{[Variable]} & \textbf{[Parametro 1]} & \textbf{[Parametro 2]} & \textbf{[Parametro 3]} & \textbf{[Estado]} \\
        \midrule
        <Fila 1> & <Valor> & <Valor> & <Valor> & <Estado> \\
        <Fila 2> & <Valor> & <Valor> & <Valor> & <Estado> \\
        \bottomrule
    \end{tabular}
\end{table}
